\documentclass[tikz,border=7pt]{standalone}
\usepackage{ctex}
\usepackage{amsmath}
\usetikzlibrary{arrows.meta,positioning,calc}

\begin{document}
\begin{tikzpicture}[
  x=1cm,y=1cm,
  >=Stealth,
  every node/.style={font=\rmfamily\footnotesize, align=center, text=black},
  title/.style={font=\bfseries\small, anchor=west},
  subtitle/.style={font=\scriptsize, text=black!58, anchor=west},
  lane/.style={font=\bfseries\scriptsize, anchor=west},
  source/.style={
    rectangle, rounded corners=2pt,
    draw=black!40, fill=white, line width=0.58pt,
    minimum width=2.15cm, minimum height=0.60cm,
    inner sep=2.5pt
  },
  mainnode/.style={
    rectangle, rounded corners=2pt,
    draw=mainGold, fill=white, line width=0.70pt,
    minimum width=2.35cm, minimum height=0.74cm,
    inner sep=3pt
  },
  rolenode/.style={
    rectangle, rounded corners=2pt,
    draw=black!45, fill=white, line width=0.58pt,
    minimum width=2.35cm, minimum height=0.72cm,
    inner sep=3pt
  },
  basenode/.style={
    rectangle, rounded corners=2pt,
    draw=black!38, fill=black!3, line width=0.50pt,
    minimum width=2.10cm, minimum height=0.56cm,
    inner sep=2.5pt
  },
  arr/.style={-Stealth, draw=black!52, line width=0.68pt, shorten >=2pt, shorten <=2pt},
  mainarr/.style={-Stealth, draw=mainGold!78!black, line width=0.80pt, shorten >=2pt, shorten <=2pt},
  rolearr/.style={-Stealth, draw=black!45, line width=0.62pt, shorten >=2pt, shorten <=2pt},
  basearr/.style={-Stealth, draw=black!42, line width=0.55pt, dashed, shorten >=2pt, shorten <=2pt},
  guide/.style={draw=black!16, line width=0.45pt}
]
  \definecolor{mainGold}{RGB}{177,129,33}
  \definecolor{kBlue}{RGB}{55,100,170}
  \definecolor{kBlueFill}{RGB}{228,237,251}
  \definecolor{vTeal}{RGB}{62,133,128}
  \definecolor{vTealFill}{RGB}{227,243,241}

  \node[title] at (0,5.16) {行为敏感度画像到预算分配};
  \node[subtitle] at (0,4.80) {主线读取 K 侧逐层画像，K/V 双侧画像用于角色预算扩展};
  \draw[guide] (0,4.56) -- (14.10,4.56);

  \node[lane, text=black!70] at (0.12,4.18) {离线校准产物};
  \node[source, draw=kBlue, fill=kBlueFill] (sigk) at (1.18,3.55)
    {K 侧读数\\$\sigma_K^{(l)}$};
  \node[source, draw=vTeal, fill=vTealFill] (sigv) at (1.18,2.55)
    {V 侧读数\\$\sigma_V^{(l)}$};

  \node[lane, text=mainGold!78!black] at (3.18,4.38) {逐层预算主线};
  \node[mainnode] (profile) at (4.20,3.55)
    {K 侧画像\\$S_{K,\alpha}$\\[-1pt]\scriptsize $\operatorname{Agg}_{\alpha}(\sigma_K^{(l)})$};
  \node[mainnode] (topk) at (7.05,3.55)
    {保护集合\\$\mathcal P_k$\\[-1pt]\scriptsize $\operatorname{TopK}(S_{K,\alpha},k)$};
  \node[mainnode] (bits) at (9.85,3.55)
    {逐层位宽\\$b^{(l)}$\\[-1pt]\scriptsize INT4 / INT8};

  \draw[mainarr] (sigk.east) -- (profile.west);
  \draw[mainarr] (profile.east) -- (topk.west);
  \draw[mainarr] (topk.east) -- (bits.west);

  \node[basenode] (heur) at (4.20,1.56)
    {Heuristic-$k$\\位置基线};
  \node[basenode] (hset) at (7.05,1.56)
    {同预算\\对照集合};
  \node[basenode] (hbits) at (9.85,1.56)
    {对照位宽\\$b_{\mathrm{pos}}^{(l)}$};
  \draw[basearr] (heur.east) -- (hset.west);
  \draw[basearr] (hset.east) -- (hbits.west);

  \node[lane, text=black!68] at (3.18,2.44) {位置启发式对照};
  \draw[guide] (3.18,2.20) -- (11.00,2.20);

  \node[lane, text=black!68] at (3.18,0.74) {K/V 角色预算扩展接口};
  \node[rolenode] (roleprofile) at (4.20,-0.02)
    {双侧画像\\$S_{K,\alpha},\,S_{V,\alpha}$};
  \node[rolenode] (roleset) at (7.05,-0.02)
    {角色保护集合\\$\mathcal P_K,\mathcal P_V$};
  \node[rolenode] (rolebits) at (9.85,-0.02)
    {角色位宽接口\\$\mathbf b^{(l)}=(b_K^{(l)},b_V^{(l)})$};

  \coordinate (sourcepair) at ($(sigk.east)!0.50!(sigv.east)+(0.24,0)$);
  \draw[rolearr] (sourcepair) to[out=-72,in=176] (roleprofile.west);
  \draw[rolearr] (roleprofile.east) -- node[above, font=\scriptsize, text=black!55] {$k_K,k_V$} (roleset.west);
  \draw[rolearr] (roleset.east) -- (rolebits.west);

  \node[font=\scriptsize, text=black!55, anchor=west] at (11.18,3.55)
    {主线实验入口};
  \draw[mainarr] (bits.east) -- (11.08,3.55);
  \node[font=\scriptsize, text=black!55, anchor=west] at (11.18,-0.02)
    {接口预留};
  \draw[rolearr] (rolebits.east) -- (11.08,-0.02);
\end{tikzpicture}
\end{document}
